\documentclass{article}


% NeurIPS 2026 style, PREPRINT mode: non-anonymous (shows author names),
% footer just says "Preprint." (no false conference claim) — right for a course report.
% (NeurIPS 2026 'final' would require a track option like [final,main] and print
%  the conference name in the footer, which we don't want here.)
\usepackage[preprint]{neurips_2026}


\usepackage[utf8]{inputenc}   % allow utf-8 input
\usepackage[T1]{fontenc}      % use 8-bit T1 fonts
\usepackage{hyperref}         % hyperlinks
\usepackage{url}              % simple URL typesetting
\usepackage{booktabs}         % professional-quality tables
\usepackage{amsfonts}         % blackboard math symbols
\usepackage{amsmath}
\usepackage{nicefrac}         % compact symbols for 1/2, etc.
\usepackage{microtype}        % microtypography
\usepackage{xcolor}
\usepackage{graphicx}
\graphicspath{{figures/}}


% ==== Working title (team can change) ====
\title{Prompt-Controllable Synthetic HAR Time Series with LLMs:\\
Capabilities and Representation Limits of SDForger on PAMAP2}


% ==== Authors (TODO: fill exact names / order / emails) ====
\author{%
  Huatai (Wichai) Pan \\
  Aalto University \\
  \texttt{firstname.lastname@aalto.fi} \\
  \And
  Teammate~2 \\ Aalto University \\
  \And
  Teammate~3 \\ Aalto University \\
  % ELEC-E7633 Project Course. Advisor: Si Zuo.
}


\begin{document}


\maketitle


\begin{abstract}
% TODO: 150-200 words. Problem (LLM-based synthetic HAR time series, SDForger on
% PAMAP2), what we did (prompt/stat conditioning, gpt2->Qwen, per-activity repair),
% and the headline finding: prompt controls amplitude but NOT cadence/frequency,
% which is a *representation* limit of the fixed linear FICA basis (empirically:
% run reconstructed through a walk-fit basis collapses to walking's cadence).
\end{abstract}


\section{Introduction}
\label{sec:intro}
% TODO: research question. Synthetic HAR data for privacy / data scarcity.
% SDForger = FICA embedding + LLM over serialized coefficients. Our theme = prompt
% design / controllability. State the contribution up front:
% (1) dual-layer controllability characterization (amplitude vs frequency);
% (2) activity-level cadence collapse; (3) linear representation ceiling + a
% nonlinear-embedding direction (validated PoC / future work).


\section{Related Work}
\label{sec:related}
% TODO: SDForger (linear FICA/FPC + LLM); time-series generation (TimeVAE,
% TimeVQVAE, TimeGAN, diffusion); TSGBench metrics; audio codec+LM analogy
% (EnCodec/AudioLM) as the nonlinear-tokenizer precedent. PAMAP2 dataset.


\section{Methods}
\label{sec:methods}
% TODO: SDForger pipeline (periodicity segmentation -> FastICA ~8 coeffs ->
% text serialization -> LLM next-token -> linear decode window=coeffs@mixing.T+mean).
% Our additions: prompt/stat conditioning; gpt2 vs Qwen2.5-1.5B (HF backend);
% per-activity clip repair. Metrics: TSGBench (MDD/ACD/SD/KD/ED/DTW); adherence
% = Spearman(requested level, realized stat); HAR downstream accuracy.


\section{Experiments}
\label{sec:experiments}
% TODO (advisor: list models tried, params tuned, ALL results):
% - gpt2 vs Qwen (explosion vs stable); multi-subject; normalization ablation.
% - Best config (Qwen + multi-subject + clip repair): train/held-out ACD.
% - Prompt controllability: max/std/range/period adherence (+ per-descriptor TSG).
% - Trimming lead-in (clean ACF-distance metric); activity discriminability.
% - Honest baselines (Gaussian-coeff, naive jitter, conditional Gaussian, NN memo).
% - Representation-ceiling PoC (run through walk-fit FICA vs AE; R^2 + period).


\section{Result Analysis}
\label{sec:analysis}
% TODO (advisor: WHY good / WHY bad, not just numbers). This is the heart.
% - Dual-layer controllability: amplitude in coefficients (controllable ~0.40);
%   frequency welded in the fixed linear basis (uncontrollable ~0). Training layer:
%   ~50-epoch sweet spot, over-training rigidifies.
% - Activity-level consequence: generated walk/run collapse to the same cadence
%   (~62/63 vs real 55/77) -> not distinguishable downstream (~chance).
% - Representation ceiling (proven): linear affine decode cannot represent an
%   out-of-basis fundamental; nonlinear is NECESSARY but not sufficient.
% - Methodology note: TSGBench ACD is unreliable for cross-model (differently
%   normalized) comparison; use scale-invariant raw-space measures.


\section{Conclusion and Future Work}
\label{sec:conclusion}
% TODO: summary of the capability/limit story. Future: nonlinear embedding
% (VQ-VAE) + conditioning to break the cadence ceiling; multichannel generation
% (SDForger's signature gap); more subjects.


% \nocite{*}  % uncomment to list all refs.bib entries even if not cited
\bibliographystyle{plainnat}
\bibliography{refs}


\end{document}
